%% ma: L12-B08-0149
%% lop: 12
%% chuong: 2
%% bai: 8
%% dang: 12.8.50
%% loai: TN4
%% muc: VD
%% dap_an: "A"
%% nguon: "CĐ2.3.pdf giáo viên gửi 27/10/2026 (M07, Lê Minh Khôi – Nguyễn Anh Khôi, Chương II Chủ đề 3 Biểu thức tọa độ của các phép toán vectơ), Câu 65"
%% kien_thuc: "Bài 8 (điểm chia đoạn thẳng, độ dài), lớp 10 (tính chất đường phân giác)"
%% trang_thai: cho-duyet
%% ly_do: "Dạng mới 12.8.50 đề xuất, chờ giáo viên duyệt"
\begin{ex}
Trong không gian $Oxyz$, cho tam giác $ABC$ có $A(1;2;-1)$, $B(2;-1;3)$, $C(-4;7;5)$. Tọa độ chân đường phân giác trong góc $B$ của tam giác $ABC$ là
\choice{\True $M\left(-\dfrac23;\dfrac{11}3;1\right)$}{$N\left(\dfrac{11}3;-2;1\right)$}{$K\left(\dfrac23;\dfrac{11}3;\dfrac13\right)$}{$V(-2;11;1)$}
\loigiai{$BA=\sqrt{26}$, $BC=\sqrt{104}=2\sqrt{26}$ nên chân phân giác $M$ thỏa $\overrightarrow{MC}=-2\overrightarrow{MA}$.

$M=\dfrac{2A+C}3=\left(-\dfrac23;\dfrac{11}3;1\right)$.}
\end{ex}
